\documentclass[11pt]{article}
\usepackage[margin=1in]{geometry}
\usepackage{amsmath,amssymb,amsthm,microtype,hyperref}
\hypersetup{hidelinks,pdfauthor={T},pdftitle={Causal Benefit and Policy Risk Under Optional Assistance}}
\newtheorem{theorem}{Theorem}
\title{Causal Benefit and Policy Risk Under Optional Assistance}
\author{T}
\date{}
\begin{document}
\maketitle
\begin{abstract}
An assistance offer can help people who accept it while harming
people who would decline. Under one-sided uptake we derive the
sharp range of outcomes from carrying out the intervention for
everyone, and use it to bound the downside of an authorized policy
change. For two voluntary offers we bound the benefit to newly
induced takers when the offer itself can affect their outcomes.
A separate result gives the possible shares of people helped and
harmed, even when average benefit is positive. Exact finite checks
reproduce the bounds. Applying them requires evidence for the
direct-effect allowances and separate authorization for each regime.
An exact allocation then balances worst-case policy regret against
an overall individual-harm allowance across groups.
\end{abstract}

\section*{Summary and scope}
\paragraph{Question.} Does a favorable trial of voluntary help justify
changing which people receive it?
\paragraph{Finding.} With one-sided uptake, the unobserved outcomes of
non-takers give a sharp interval for mandatory execution. With two
voluntary offers, bounded direct effects give a different sharp interval
for newly induced takers. Its two outcome marginals give a sharp
joint region for the shares who benefit and who are harmed.
\paragraph{Evidence and scope.} The bounds are proved and checked
against exact synthetic oracles. An operator can report benefit
sensitivity alongside uptake before expanding an offer. The required
direct-effect bounds and individual consent come from outside the trial.

\section{What a randomized offer identifies}
A scheduling assistant can offer a change, let a person decline,
and observe what happens. A positive result for the offer does not
determine what would happen if the assistant carried out the change
for everyone. The distinction is consequential when a service uses
voluntary trial data to change an automated default.

Let $Z$ be a randomized offer, $D(z)$ indicate uptake under offer status $z$, and $Y(d)$ be the outcome under treatment status $d$. Assume no interference, $D(1)\ge D(0)$, and that the offer affects the outcome only through uptake. Always-takers and never-takers contribute zero to the intent-to-treat contrast because their treatment does not change. Compliers contribute $Y(1)-Y(0)$. Thus, for nonzero uptake contrast,
\[
\frac{\mathbb E[Y\mid Z=1]-\mathbb E[Y\mid Z=0]}
     {\mathbb E[D\mid Z=1]-\mathbb E[D\mid Z=0]}
 =\mathbb E[Y(1)-Y(0)\mid D(1)>D(0)].
\]
This is the local average treatment effect of Imbens and
Angrist \cite{imbens}, in the instrumental-variable framework
reviewed by Pearl \cite{pearl}. The distinction matters for
learned assistance policies because data from voluntary users
may favor an action that other users reject.

\paragraph{A reversal in a finite population.}
In a finite population of 100 simulated recipients, 30 compliers gain two units and 70 never-takers lose two if treated. The uptake contrast is $0.30$ and the offer's outcome contrast is $0.60$, giving the ratio $2$. Forcing treatment on everyone produces total change $60-140=-80$. Both computations use exact integer or rational arithmetic.


\section{What a favorable offer leaves unresolved}
Consider one-sided uptake, so $D(0)=0$ and $D(1)=1$ for compliers
only. Let $c$ be the complier share. Outcomes lie in $[0,1]$.
Write $\mu_{d,C}$ and $\mu_{d,N}$ for mean outcomes under
treatment $d$ in compliers and never-takers. The offer arm
identifies $c$, $\mu_{1,C}$, and $\mu_{0,N}$ when the corresponding
strata have positive mass. The control arm identifies the population
untreated mean $W$. This in turn identifies $\mu_{0,C}$ for $c>0$.
The policy that offers assistance has mean
\[
O=c\mu_{1,C}+(1-c)\mu_{0,N}.
\]
Mandatory execution would yield
$F=c\mu_{1,C}+(1-c)\mu_{1,N}$, where $\mu_{1,N}$ is unobserved.
Its sharp interval is therefore
\[
[L,U]=[c\mu_{1,C},\ c\mu_{1,C}+1-c].
\]
The endpoint worlds set the missing mean to zero or one without
altering the observed arm distributions. Intermediate means are
possible by mixing those worlds. The limiting case $c=0$ uses
$c\mu_{1,C}=0$ and the case $c=1$ collapses the interval.
This is the one-sided special case of the partial-identification
problem studied by Balke and Pearl \cite{balke}.

\section{A risk limit on policy experimentation}
Choose the better of the two identified regimes and denote its mean
by $S=\max(W,O)$. Consider a policy that uses the mandatory regime
with probability $\theta$ and the identified regime otherwise,
where $0\le\theta\le1$. This is an analysis of authorized regime
choices. A standing refusal or missing authority excludes the
mandatory regime before numerical optimization.

Suppose $L<S<U$. Define $a=S-L$ and $b=U-S$. The comparator
knows $F$ and chooses whichever regime has the higher mean.
For $F\le S$, regret is $\theta(S-F)$ and increases as $F$
decreases. For $F\ge S$, regret is $(1-\theta)(F-S)$ and
increases as $F$ increases. The worst-case regret is therefore
\[
R(\theta)=\max\{\theta a,(1-\theta)b\}.
\]
A separate constraint limits the worst-case expected loss against
$S$ to $\delta\ge0$, which gives $\theta a\le\delta$.
The unconstrained minimizer is the intersection of the two linear
pieces. The objective decreases up to that point and increases
after it. The constrained optimizer is
\[
\theta^*=\min\!\left\{\frac{b}{a+b},\frac{\delta}{a}\right\}.
\]
If $L\ge S$ choose the mandatory regime whenever authorized.
If $U\le S$ choose the identified regime. The formula applies the
minimax-regret criterion for treatment decisions \cite{manski}
with an additional downside constraint. It controls an expected
policy contrast over the modeled population. In particular it
cannot by itself grant individual permission to act.

\section{Exact evaluation of the decision rule}
Take $c=2/5$, $\mu_{0,C}=1/5$, $\mu_{1,C}=4/5$, and
$\mu_{0,N}=3/5$. Offering improves the mean from $11/25$ to
$17/25$. Mandatory execution can produce any mean in $[8/25,23/25]$.
Without the additional risk limit, $\theta^*=2/5$. With
$\delta=1/20$, the feasible optimizer is $5/36$, with worst
regret $31/150$. If mandatory execution is excluded by consent,
the authorized policy set uses $\theta=0$ and worst regret
against the unrestricted oracle is $6/25$.

The checker evaluated 625 parameter settings and five loss caps
per setting. The closed form attained no greater regret than any
admissible member of a 17-point grid in all 3,125 cases.
The risk cap reduced the unconstrained choice in 271 cases.
Ninety parameter settings had strictly favorable offering and
an interval for mandatory execution that straddled the offer's
mean. These are finite design counts. The mathematical derivation
covers continuous mixtures and the experiment checks its implementation.

\section{Two voluntary offers with direct effects}
An offer can change distress even when the recipient does not
accept it. Instead of assuming exact exclusion, compare a weak
and a strong authorized offer, labeled $z=0,1$. Keep randomized
assignment, noninterference, and monotone uptake. The three
uptake strata are stable takers $A$, induced takers $C$, and
stable non-takers $N$, with shares $a,c,n$. They are identified
from the two uptake rates, and assume $c>0$. Write $Y_z(d)$
for a binary outcome under offer $z$ and treatment $d$.

Use population-weighted success masses. Let $A_z$ be the
success mass for $A$ under $Y_z(1)$ and $N_z$ the mass for $N$
under $Y_z(0)$. Let $C_{zd}$ be the success mass for $C$ under
$Y_z(d)$. The weak treated cell identifies $A_0$, and the
strong untreated cell identifies $N_1$. The other two cells
identify $h_{11}=A_1+C_{11}$ and $h_{00}=N_0+C_{00}$.
For example $h_{11}=\Pr(Y=1,D=1\mid Z=1)$.

Supply three sensitivity parameters, in population-mass units:
\[
 |A_1-A_0|\le \epsilon_A,\qquad
 |N_1-N_0|\le \epsilon_N,\qquad
 |C_{10}-C_{00}|\le \epsilon_C.
\]
These are assumptions about direct offer effects at the same
treatment status. The trial alone does not identify them.
Define
\begin{align*}
 I_A &= [A_0-\epsilon_A,A_0+\epsilon_A]\cap[0,a],\\
 I_N &= [N_1-\epsilon_N,N_1+\epsilon_N]\cap[0,n],\\
 I_1 &= (h_{11}-I_A)\cap[0,c],\\
 I_0 &= (h_{00}-I_N)\cap[0,c],\\
 J_0 &= (I_0+[-\epsilon_C,\epsilon_C])\cap[0,c].
\end{align*}
Here sums and differences range over all points in the intervals.

\begin{theorem}
If either $I_1$ or $I_0$ is empty, the observed cells are
incompatible with the stated bounds. Otherwise the sharp
identified interval for the treatment benefit among induced
takers under the strong offer is
\[
 \mathbb E[Y_1(1)-Y_1(0)\mid C]
 \ \in\ \frac{1}{c}(I_1-J_0).
\]
\end{theorem}
\begin{proof}
The strong treated mixture gives $C_{11}=h_{11}-A_1$,
and the weak untreated mixture gives $C_{00}=h_{00}-N_0$.
This yields $I_1$ and $I_0$ by the observed totals and the
direct-effect bounds. The constraint on the untreated induced
stratum expands $I_0$ to $J_0$. The choices of $C_{11}$
and $C_{10}$ are independent as feasibility constraints:
they involve disjoint unknown masses and can be combined
in one joint potential-outcome law with the specified
marginals. Thus every pair of points in $I_1$ and $J_0$
is feasible, including the pairs that attain the endpoints.
\end{proof}

For a trial with $(a,c,n)=(1/4,1/4,1/2)$, suppose
$(A_0,h_{11},h_{00},N_1)=(0.15,0.35,0.35,0.25)$.
The uptake contrast is $1/4$ and the outcome contrast is $0.10$.
Under exact exclusion their ratio gives benefit $0.40$.
Allowing same-treatment mean shifts of $0.10$ for $A$ and $N$
and $0.05$ for untreated $C$ gives benefit in $[0.05,0.75]$.
Increasing only the allowable mean shift for $N$ to $0.20$
widens the interval to $[-0.15,0.90]$.

An exact integer-grid audit enumerated latent success counts
and independently compared both endpoints or infeasibility for
56,700 observed count and sensitivity settings. It found 4,299
settings with a positive conventional Wald contrast and a
strictly sign-ambiguous interval. That count describes the
declared grid. The sharpness argument works over continuous
population masses and does not rely on the grid.

\section{A positive mean can conceal individual harm}
For the strong offer and induced takers, write $[p_-,p_+]=I_1/c$
for the treated success probability and $[q_-,q_+]=J_0/c$
for the untreated one. A person benefits when the treated outcome
is one and the untreated outcome is zero, and is harmed in the
reverse case. Let $b$ and $h$ be the corresponding population shares
within the induced stratum. The mean treatment effect is $b-h$.
It leaves the two shares underdetermined.

\begin{theorem}
With unrestricted dependence of the two potential outcomes, the
sharp region for $(b,h)$ is
\begin{align*}
 0\le b &\le \min(p_+,1-q_-),\\
 0\le h &\le \min(q_+,1-p_-),\\
 b+h &\le 1,\\
 p_--q_+ &\le b-h\le p_+-q_- .
\end{align*}
The sharp marginal lower bounds are $\max(0,p_--q_+)$ for $b$
and $\max(0,q_--p_+)$ for $h$. The first two inequalities give
their respective sharp upper bounds.
\end{theorem}
\begin{proof}
Let $u$ be the share with both potential outcomes equal to one.
The four joint masses are $b,h,u,1-b-h-u$. With $b,h\ge0$,
their nonnegativity and the marginal constraints hold exactly when
\[
 [0,1-b-h]\cap[p_--b,p_+-b]\cap[q_--h,q_+-h]
\]
is nonempty. Requiring each lower endpoint to be at most each
upper endpoint gives the displayed inequalities. Any allowed
$u$ constructs a joint law, which can be coupled with the
remaining binary potential outcomes to match their feasible
marginals in the preceding theorem.

For fixed success marginals $p,q$, the smallest and largest
possible benefit shares are $\max(0,p-q)$ and $\min(p,1-q)$.
This follows by choosing the largest and smallest feasible
common-success share $u$, respectively. Minimizing the former
over the marginal intervals uses $p_-,q_+$. Maximizing the
latter uses $p_+,q_-$. Interchanging the roles of treatment
and no treatment proves the harm bounds.
\end{proof}

In the trial with the narrower direct-effect allowances, the
induced group's mean treatment effect lies in $[0.05,0.75]$,
its benefited share in $[0.05,0.85]$ and its harmed share in
$[0,0.30]$. A joint law with $b=0.35$, $h=0.30$ and $u=0.35$
matches the observed cells and attains the lowest mean effect.
The largest benefit and harm bounds cannot occur together
because their sum would exceed one. Thus an operator with a
ten-percent ceiling on individual harm cannot certify that
ceiling from the positive mean alone.

The verifier enumerated joint outcome tables on populations of
one to eight induced takers. The region and all four coordinate
extrema matched direct enumeration for 4,916 pairs of marginal
intervals and 308,846 candidate benefit and harm pairs. Of
those interval pairs, 560 permit positive harm even though
every compatible law has positive mean benefit. This is a finite
algebra check. Trial sampling uncertainty and the direct-effect
assumptions still need separate analysis in an application.

\section{Allocating authorized help under a harm allowance}
Suppose a service has pre-decision groups with shares $w_g$ that
sum to one. For each group, assume justified treatment-effect bounds
$[\ell_g,u_g]$ and a sharp upper bound $\bar h_g$ on the share of
people whose outcome would worsen under treatment. For binary
outcomes, the preceding joint region supplies these quantities.
Take the joint uncertainty across groups to be the product of
their admissible laws. Compare treatment with a fixed reference
under the same offer. Within each group the action is authorized,
and the service assigns treatment with probability $\theta_g$
independently of each person's potential outcomes.

Let $a_g=[-\ell_g]_+$ and $b_g=[u_g]_+$, where $[x]_+=\max(x,0)$.
An oracle that knows the mean effect can choose to treat a whole
group or leave it at the reference. The worst regret against this
oracle and the worst share of people harmed are, respectively,
\[
 \sum_g w_g\max\{a_g\theta_g,b_g(1-\theta_g)\},
 \qquad \sum_g w_g\bar h_g\theta_g.
\]
For the first expression, consider the extreme group means.
For the second, independence of assignment gives a harm share
of $\theta_g h_g$ in each group. Product uncertainty makes both
bounds sharp. The two suprema can occur in different worlds.

\begin{theorem}
To minimize worst-case regret subject to an overall harm allowance
$\eta$, let $x_g=b_g/(a_g+b_g)$ when the denominator is positive,
and set $x_g=0$ otherwise. Use $\theta_g=x_g$ for groups with
$\bar h_g=0$. For the others, assign treatment probability up to
$x_g$ in descending order of $b_g/\bar h_g$, stopping when
$\sum_g w_g\bar h_g\theta_g=\eta$ or all group limits are reached.
\end{theorem}
\begin{proof}
Moving any $\theta_g>x_g$ down to $x_g$ weakly lowers both the
regret and its harm cost. On $[0,x_g]$, increasing $\theta_g$
saves $w_g b_g$ units of regret per unit of probability and costs
$w_g\bar h_g$ units of the harm allowance. There is no cost when
$\bar h_g=0$. For positive cost, moving an equal amount of harm
allowance from a lower-ratio group to a higher-ratio group with
room left cannot increase regret. Repeating the exchange gives
the stated allocation, including ties and a partial final group.
\end{proof}

For two equally sized groups, take treated and untreated success
intervals $[0.40,0.80]$ and $[0.20,0.40]$ in the first group,
and $[0.90,0.90]$ and $[0.60,0.80]$ in the second. The upper mean
improvements are $0.60$ and $0.30$. The maximum harmed shares
are $0.40$ and $0.10$. At overall allowance $0.08$, serving the
second group first gives treatment rates $(0.15,1)$ and regret
$0.255$. Serving the larger upper mean first gives rates $(0.40,0)$
and regret $0.330$. Both meet the same harm allowance.

The verifier reconstructed 384 group extrema from joint outcome
tables and matched 640 allocations against a vertex oracle for the
piecewise-linear objective. In an application, use group bounds
for the actual authorized population. Principal strata from a trial
can support a sensitivity analysis, but observed group membership
needs its own evidence. The overall allowance can be supplemented
by separate harm ceilings in groups requiring additional protection.

\section{Application and next test}
Train the policy to distinguish offering, accepting, and executing. During a controlled evaluation of an accessibility assistant, log opt-out friction and downstream outcomes separately. If receiving an offer itself changes distress, exclusion fails and the randomized effect of the offer remains the defensible estimand. Cross-user coordination can also violate noninterference. A follow-up should test a low-friction offer against a matched control and report both uptake and outcome contrasts. A negative offer effect would directly challenge the proposed deployment even with favorable outcomes among acceptors.

\paragraph{Review status.} Independent manuscript review is pending.

\begin{thebibliography}{9}
\bibitem{pearl} Judea Pearl. \emph{Causal inference in statistics: An overview}. 2009. \url{https://ftp.cs.ucla.edu/pub/stat_ser/r350.pdf}.
\bibitem{imbens} Guido W. Imbens and Joshua D. Angrist.
Identification and Estimation of Local Average Treatment Effects.
\emph{Econometrica} 62(2), 467--475, 1994.
\url{https://raw.githubusercontent.com/Mixtape-Sessions/Instrumental-Variables/main/Readings/Lecture3/Angrist_Imbens_1994.pdf}.
\bibitem{balke} Alexander Balke and Judea Pearl.
\emph{Bounds on Treatment Effects From Studies With Imperfect Compliance}.
Journal of the American Statistical Association 92(439), 1997.
\url{https://ftp.cs.ucla.edu/pub/stat_ser/r199-jasa.pdf}.
\bibitem{manski} Charles F. Manski.
\emph{Statistical Treatment Rules for Heterogeneous Populations}.
\emph{Econometrica} 72(4), 1221--1246, 2004.
\url{https://www.cemmap.ac.uk/wp-content/legacy/resources/files/manski09_5.pdf}.
\end{thebibliography}
\end{document}
